\section{Introduction}
\label{sec:intro}

Ask a video assistant ``how long did I use the paper cutter?'' and a useful answer gives a duration and the interval
it comes from. Such \emph{grounded} questions are central to egocentric assistants, video search and step-by-step
guidance, and they require knowing when things happened across the whole window.

Video LLMs pay for every frame. Qwen3-VL~\cite{qwen3vl} spends about 330 visual tokens per frame, so 64 frames
cost about 21{,}000 tokens. Pruning has therefore become the default, and many works report that 10--25\% of the
visual tokens preserve multiple-choice accuracy~\cite{fastvid,vidcom2,flashvid,semvid,prunevid,dycoke}. We ask what
pruning does to \emph{temporal grounding}. In short, \textbf{pruning for QA keeps what is salient, but grounding needs
what is everywhere.} A QA answer can be read off a few informative frames. An interval needs a timeline with an anchor
in every part of the window.

\paragraph{Observation 1: QA pruners lose the anchors.}
Video LLMs such as Qwen3-VL read time from timestamp tokens interleaved with the visual tokens. A frame that loses all
its visual tokens loses its timestamp too, or keeps it with nothing to show. QA pruners drop frames that look
redundant, unremarkable or unrelated to the question, which are exactly the frames a timeline needs. All four recent
pruners we evaluate fall below random token selection at the same budget on S-EMBER, QVHighlights and Charades-STA
(\cref{sec:pruners}). Random selection leaves a few tokens on every frame, and one published pruner gains $5.3$ mIoU
just by ranking tokens within each frame instead of globally (\cref{sec:coverage-loss}).

\paragraph{Observation 2: a prior fills the silence.}
Grounded QA invites reasoning-style answers that list observations and their times before the interval.
Asked this way without video, the model still gives confident intervals, and 90\% of them start in the first
5\% of the window. Pruned memories share the habit: 30--56\% of their intervals start there, against 16\% of the gold
intervals, and removing the evidence frames raises this share from 37\% to 51\% (\cref{sec:prior}). The prior is
not tied to any pruner. It is what the decoder writes when the memory is silent.

\paragraph{Observation 3: the format decides which failure dominates.}
On the same QVHighlights queries and memory, asking for a timeline of observations instead of a single interval costs
about 30 mIoU (\cref{sec:transfer}). With one-line answers, question-driven token selection helps and contrastive
decoding does nothing. With timeline answers, the prior writes many of the times, so contrasting it away helps unless
the contrast breaks the format. We therefore treat the answer format as an experimental variable.

\medskip
We propose \method (BOLT $+$ Coverage $+$ Contrastive decoding), a training-free method with one component per
failure and an optional BOLT front end that chooses which frames to encode (\cref{fig:overview}).
\begin{itemize}
  \item \textbf{Coverage} (\cref{sec:coverage}) keeps the anchors. It prunes the KV cache layer by layer, gives each
  frame an equal share of the budget, and keeps within each frame the tokens the question attends to most. The quota
  guarantees every time anchor, and the question decides what each anchor shows. In the settings we tested, a global
  top-$k$ over the same attention does as well, so the question score drives the gain and the quota is a safeguard
  against scores collapsing onto a few frames.
  \item \textbf{Prior-contrastive decoding (PCD)} (\cref{sec:contrast}) removes the prior. It chooses each time value
  by contrasting the model's prediction with its prediction from a video-free memory, as in contextual calibration
  and context-aware decoding~\cite{calibratebeforeuse,cad}, with a strength that fades where the model is confident.
\end{itemize}

Neither operator is new on its own. Per-frame budgets appear in VidCom\textsuperscript{2}, KTV and
AdaTP~\cite{vidcom2,ktv,adatp}, question-attention KV eviction in SnapKV and PruneVid~\cite{snapkv,prunevid}, and
contrast against a content-free or degraded input in contrastive decoding, classifier-free guidance and
VCD~\cite{contrastivedecoding,cfglm,cad,vcd,m3id}. We contribute an account of \emph{why} grounding needs each, when
each matters, and how to combine them:
\begin{enumerate}
  \item \textbf{A diagnosis of grounding under compression} (\cref{sec:diagnosis}): frame-emptying pruners and,
  shown causally by evidence removal, a video-free prior in reasoning-style answers.
  \item \textbf{Coverage} (\cref{sec:coverage}). With the one-line prompt on the full Charades-STA and QVHighlights
  sets, for Qwen3-VL-8B and the grounding-tuned TimeLens-8B, it matches the uncompressed model with 24\% of the visual
  tokens and beats random selection by $+0.5$ to $+1.6$ mIoU (\cref{sec:transfer}). On NExT-GQA at 10\% of the tokens
  it keeps the uncompressed accuracy, while random selection loses about four points. Against a 1\,fps input of
  similar budget, most of the gain comes from the BOLT front end.
  \item \textbf{PCD and an answer-format study} (\cref{sec:contrast,sec:transfer}). PCD helps reasoning-style
  S-EMBER answers on all eight memories we test and on both models. With 10\% of the tokens, \method improves on the
  uncompressed model by $+11.2$ and $+10.7$ mIoU under the same prompt, and by $+5.5$ under the official S-EMBER
  prompt, where the smaller accuracy gain ($+2.9$) comes mostly from BOLT. PCD does nothing for one-line answers and
  lowers NExT-GQA grounding where it disturbs the format (TimeLens-8B timeline answers, minutes:seconds under the
  official prompt).
  \item \textbf{Reproducible evaluation} (\cref{sec:setup,sec:transfer}). We reproduce published TimeLens-Bench
  numbers within 0.3 mIoU, run no-video and shuffled-frame or shuffled-timestamp sanity checks, and report negative
  results, including the limits of Coverage on short, densely sampled videos.
\end{enumerate}
